Production large language models are reliably sycophantic: they affirm user beliefs even when those beliefs are wrong, which makes them unsuitable as collaborators in tasks where the user is mistaken. A common intuition holds that giving an assistant a stronger persona should reduce this failure mode, but recent work \citep{han2025personality} reports that persona injection moves Big-Five \emph{self-reports} (\(\beta=3.95\)) while barely moving \emph{behavior} (\(\beta=0.03\), \(p=0.67\) on sycophancy), suggesting that LLM personas are linguistic surface, not action control. We test whether this null generalizes to \emph{belief-based} personas that name explicit epistemic commitments rather than demographic tags or Big-Five traits. We compare five system-prompt conditions on a ``distinct-identity spectrum'' (no persona, ``you are a helpful assistant,'' demographic, constitutional, belief-based) on two production OpenAI models across three probes: ``are you sure?'' sycophancy, productive correction of stated misconceptions, and persona stickiness measured by response-embedding distance to a held-out style reference. The belief persona reduces sycophantic agreement with confidently-wrong users by \(18.75\)--\(21.25\) percentage points on both \gptone (\(p=0.0018\)) and \gptmini (\(p=0.007\)), a result that survives Holm correction. It is the only condition whose responses cluster near the held-out persona style (Cohen's \(d \geq 1.27\) vs no-persona). Demographic and shallow ``helpful assistant'' personas show no significant behavioral effect, replicating Han et al. for those persona families. Our findings refute a strong form of the Personality Illusion: when the persona names what counts as evidence, it controls action.
